\documentclass[11pt]{article}

\usepackage{amsmath, amssymb}
\usepackage{geometry}
\usepackage{listings}
\usepackage{xcolor}
\usepackage{hyperref}
\geometry{margin=1in}

\definecolor{codebg}{rgb}{0.96,0.96,0.97}
\definecolor{codeframe}{rgb}{0.8,0.8,0.82}

\lstset{
  basicstyle=\ttfamily\small,
  backgroundcolor=\color{codebg},
  frame=single,
  rulecolor=\color{codeframe},
  breaklines=true,
  columns=fullflexible,
  showstringspaces=false,
  xleftmargin=1em,
  xrightmargin=1em,
  aboveskip=1em,
  belowskip=1em,
}

\title{Blueboard Phase 1, explained}
\author{}
\date{}

\begin{document}
\maketitle

\section*{What this document is}
Phase 1 (\texttt{PLAN.md}, tag \texttt{v0.1-window-and-camera}) took the project from an
empty cargo workspace to an actual window: a GPU-cleared screen with a reference grid you
can pan by dragging and zoom by scrolling, clamped to the brief's 1\%--200\% range. This
note walks through \emph{what was written, in what order, and why} — including the Rust
language features used along the way, for a reader who knows roughly what the pieces
should do but not yet how to spell that in Rust. It is a companion to the code, not a
replacement for reading it: file and line references point at the real source.

\tableofcontents

\section{The shape of the problem}

Three things had to exist before anything would show on screen:
\begin{enumerate}
  \item A \textbf{window} — an OS-level surface to draw into and receive input from.
  \item A \textbf{GPU pipeline} — wgpu's way of saying ``here is a shader, here is how
    vertices are laid out, here is what to draw.''
  \item A \textbf{camera} — the pan/zoom state that turns ``board space'' (where strokes
    will eventually live) into ``screen space'' (where pixels are).
\end{enumerate}
These map directly onto the three-crate split decided in Phase 0: the camera went into
\texttt{blueboard-core} (pure math, no GPU, no windowing), the pipeline into
\texttt{blueboard-render} (wgpu, nothing about windows or input), and the event loop that
ties them together into \texttt{blueboard-app} (\texttt{winit}, owns the window, owns
input handling). Each section below follows that same order: core, then render, then app
— which is also the order the code was written in, since render depends on core's
\texttt{Camera} type, and app depends on both.

\section{\texttt{blueboard-core}: the camera}

File: \texttt{crates/blueboard-core/src/camera.rs}.

\subsection{Why this lives in \texttt{core}, not \texttt{render}}
The camera is just two numbers and some arithmetic: a pan offset and a zoom factor, plus
the formulas to convert between screen pixels and board coordinates. None of that touches
the GPU. Keeping it in \texttt{core} means it can be unit-tested without ever opening a
window or talking to a graphics driver — which is exactly how it was tested (see
\S\ref{sec:testing}).

\subsection{The struct}
\begin{lstlisting}
pub const MIN_ZOOM: f64 = 0.01;
pub const MAX_ZOOM: f64 = 2.0;

#[derive(Debug, Clone, Copy)]
pub struct Camera {
    pub pan: DVec2,
    pub zoom: f64,
}
\end{lstlisting}
A few things to unpack for a reader new to Rust:
\begin{itemize}
  \item \texttt{pub const} declares a compile-time constant, public outside this module.
    \texttt{MIN\_ZOOM}/\texttt{MAX\_ZOOM} encode the brief's 1\%--200\% requirement as
    actual named values instead of magic numbers scattered through the code.
  \item \texttt{\#[derive(Debug, Clone, Copy)]} is an \emph{attribute} that asks the
    compiler to auto-generate three trait implementations: \texttt{Debug} (so
    \texttt{\{:?\}} can print a \texttt{Camera} for debugging), and \texttt{Clone}+
    \texttt{Copy} (so a \texttt{Camera} value can be duplicated by simple bitwise copy —
    cheap, since it's just two numbers — rather than requiring an explicit
    \texttt{.clone()} call or being \emph{moved} out from under its owner, which is Rust's
    default for most types).
  \item \texttt{DVec2} is \texttt{glam}'s two-component \texttt{f64} vector type —
    \texttt{glam} is the math library chosen in Phase 0's research for exactly this: a
    well-tested, fast vector/matrix type instead of hand-rolling one.
  \item Every field is \texttt{pub}: there's no invariant between \texttt{pan} and
    \texttt{zoom} that setting them directly could violate, so there's no reason to hide
    them behind getter/setter methods (a very un-Rust pattern anyway).
\end{itemize}

\subsection{Panning}
\begin{lstlisting}
pub fn pan_by_screen_delta(&mut self, screen_delta: DVec2) {
    self.pan -= screen_delta / self.zoom;
}
\end{lstlisting}
\texttt{\&mut self} is Rust's way of saying ``this method needs to mutate the
\texttt{Camera} it's called on.'' Compare to \texttt{\&self} (read-only borrow, used
below) and plain \texttt{self} (takes ownership, consumes the value — not used here since
we want to keep using the camera afterwards). The borrow checker enforces that only one
\texttt{\&mut} borrow of a value can exist at a time, which is what prevents, e.g., two
parts of the event loop from mutating the camera simultaneously and corrupting it.

The math: dragging the mouse by \texttt{screen\_delta} pixels should move the board content
by that many \emph{screen} pixels, which is \texttt{screen\_delta / zoom} board units (at
2x zoom, 10 screen pixels of drag is only 5 board units). Subtracting (rather than adding)
is because \texttt{pan} is defined as ``the board point at the viewport center'' — dragging
the canvas right moves that center point left in board space.

\subsection{Zooming, anchored on the cursor}
\begin{lstlisting}
pub fn zoom_by(&mut self, factor: f64, anchor_screen: DVec2, viewport: DVec2) {
    let anchor_board_before = self.screen_to_board(anchor_screen, viewport);
    self.zoom = (self.zoom * factor).clamp(MIN_ZOOM, MAX_ZOOM);
    let anchor_board_after = self.screen_to_board(anchor_screen, viewport);
    self.pan += anchor_board_before - anchor_board_after;
}
\end{lstlisting}
This is the ``zoom towards the cursor'' behavior every decent editor has: find which board
point is currently under the cursor, change the zoom, then nudge \texttt{pan} so that same
board point is \emph{still} under the cursor. \texttt{f64::clamp} is a standard-library
method that restricts a value to a range in one call — this is where \texttt{MIN\_ZOOM}/
\texttt{MAX\_ZOOM} actually get enforced, so no code path can push zoom outside the brief's
range no matter how it's called.

\subsection{Coordinate conversion and the view matrix}
\begin{lstlisting}
pub fn screen_to_board(&self, screen: DVec2, viewport: DVec2) -> DVec2 {
    (screen - viewport * 0.5) / self.zoom + self.pan
}

pub fn view_matrix(&self, viewport: Vec2) -> Mat4 {
    let zoom = self.zoom as f32;
    let pan = self.pan.as_vec2();
    let scale = Mat4::from_scale(Vec3::new(
        2.0 * zoom / viewport.x,
        -2.0 * zoom / viewport.y,
        1.0,
    ));
    let translate = Mat4::from_translation(Vec3::new(-pan.x, -pan.y, 0.0));
    scale * translate
}
\end{lstlisting}
\texttt{screen\_to\_board} inverts the view transform: screen pixels, relative to the
viewport center, scaled down by zoom, then shifted by the pan offset. \texttt{view\_matrix}
builds the matrix the GPU actually uses, as a composition of two 4x4 matrices: translate by
\texttt{-pan} first, then scale into \emph{clip space} (wgpu's $[-1, 1]$ coordinate range,
with Y flipped since board space is Y-down like screen space, but clip space is Y-up — the
\texttt{-2.0 * zoom / viewport.y} negative sign is doing that flip). Matrix multiplication
with column vectors applies the \emph{rightmost} operation first, so \texttt{scale *
translate} means ``translate, then scale,'' matching the math worked out in
\texttt{planning/math.tex}. Note the explicit \texttt{as f32} cast: the camera itself
stores \texttt{f64} (so board coordinates keep precision far from the origin), but GPU math
is \texttt{f32} — the cast is the one place that precision gets intentionally narrowed, right
before it reaches the GPU.

\subsection{Tests}
\label{sec:testing}
\begin{lstlisting}
#[cfg(test)]
mod tests {
    use super::*;

    #[test]
    fn zoom_clamps_to_brief_range() { /* ... */ }

    #[test]
    fn zoom_keeps_anchor_point_fixed_on_screen() { /* ... */ }
}
\end{lstlisting}
\texttt{\#[cfg(test)]} tells the compiler to only include this module when running
\texttt{cargo test} — it doesn't bloat the real binary. \texttt{mod tests \{ ... \}}
declares a nested module; \texttt{use super::*;} pulls in everything from the parent module
(\texttt{Camera}, \texttt{MIN\_ZOOM}, etc.) so the tests can use them unqualified. Each
\texttt{\#[test]}-annotated function is run independently by the test harness. These are
exactly the kind of tests that justify keeping \texttt{Camera} in \texttt{core}: no window,
no GPU, just arithmetic — they run in milliseconds.

\section{\texttt{blueboard-render}: the wgpu pipeline}

File: \texttt{crates/blueboard-render/src/lib.rs}, plus the shader at
\texttt{src/shader.wgsl}.

\subsection{The GPU pipeline, conceptually}
wgpu follows the WebGPU API fairly closely. The objects involved, in the order they're
created:
\begin{description}
  \item[\texttt{Instance}] the entry point into the GPU API — picks a backend (Vulkan on
    this Linux machine).
  \item[\texttt{Surface}] the drawable target tied to the OS window.
  \item[\texttt{Adapter}] a handle to a physical GPU.
  \item[\texttt{Device} and \texttt{Queue}] the logical device used to create GPU
    resources, and the queue used to submit work to it.
  \item[\texttt{ShaderModule}] compiled WGSL (WebGPU's shading language) code.
  \item[\texttt{RenderPipeline}] ties a vertex shader, a fragment shader, a vertex buffer
    layout, and other fixed-function state (blending, primitive topology) into one
    reusable object.
  \item[\texttt{Buffer}] raw GPU memory — used here both for vertex data (the grid lines)
    and for the uniform (the camera's view matrix, updated every frame).
  \item[\texttt{BindGroup}] the mechanism by which a buffer (or texture, or sampler) is
    actually attached to a pipeline at draw time — declared via a
    \texttt{BindGroupLayout} that describes \emph{what} will be bound, satisfied by a
    \texttt{BindGroup} that provides the actual resource.
\end{description}
\texttt{Renderer::new} in the code builds exactly this chain, once, at startup.
\texttt{Renderer::render} runs every frame: write the current camera matrix into the
uniform buffer, acquire a texture to draw into, record drawing commands into a
\texttt{CommandEncoder}, submit them, and present the result.

\subsection{Async setup and \texttt{pollster}}
\begin{lstlisting}
pub async fn new<W>(window: Arc<W>, size: (u32, u32)) -> Self
where
    W: HasWindowHandle + HasDisplayHandle + Send + Sync + 'static,
{ /* ... */ }
\end{lstlisting}
Requesting a GPU adapter and device are asynchronous operations in wgpu (they may involve
waiting on the driver). Rust's \texttt{async fn} doesn't run on its own — it returns a
\emph{future} that has to be driven to completion by an executor. Rather than pull in a
full async runtime like \texttt{tokio} for two awaited calls, the app crate uses
\texttt{pollster::block\_on(...)}, a tiny crate whose entire job is ``run this future to
completion on the current thread, synchronously.'' That's a deliberate scope choice: Phase
1 doesn't need concurrent I/O, so the simplest thing that works was used instead of the
heavier standard tool.

\subsection{Generic over the window type}
The \texttt{where W: HasWindowHandle + HasDisplayHandle + Send + Sync + 'static} clause is
a \emph{trait bound}: instead of \texttt{Renderer::new} taking a concrete
\texttt{winit::window::Window}, it takes \emph{any} type \texttt{W} that implements the two
\texttt{raw-window-handle} traits wgpu needs to create a surface, plus \texttt{Send + Sync}
(safe to share across threads) and \texttt{'static} (doesn't borrow anything with a limited
lifetime). This is what keeps \texttt{blueboard-render} from depending on \texttt{winit} at
all — per \texttt{PLAN.md}'s architecture section, render ``knows nothing about UI,'' and
this is the concrete mechanism that makes that true while still letting \texttt{winit}'s
\texttt{Window} be passed in from the app crate.

\texttt{Arc<W>} (atomically-reference-counted pointer) is needed because the window has to
be kept alive for as long as the \texttt{Surface} borrows from it, \emph{and} the app crate
separately keeps its own handle to the same window to call things like
\texttt{request\_redraw()}. \texttt{Arc} allows both owners to share the same allocation
safely; \texttt{window.clone()} elsewhere in the app crate just bumps a reference count, it
doesn't copy the window.

\subsection{Vertex data and \texttt{bytemuck}}
\begin{lstlisting}
#[repr(C)]
#[derive(Clone, Copy, bytemuck::Pod, bytemuck::Zeroable)]
struct Vertex {
    position: [f32; 2],
}
\end{lstlisting}
GPU buffers are raw bytes with an agreed-upon layout — there's no type information on the
GPU side. \texttt{\#[repr(C)]} pins the struct's memory layout to C conventions (no field
reordering, which Rust is otherwise free to do), so it matches exactly what the shader
expects. \texttt{bytemuck::Pod} (``plain old data'') and \texttt{Zeroable} are marker
traits, checked at compile time, that certify a type is safe to reinterpret as a byte
slice — which is what \texttt{bytemuck::cast\_slice(\&grid\_vertices)} does when uploading
the grid to a GPU buffer. Without those guarantees, reinterpreting arbitrary structs as
bytes would be undefined behavior (e.g. if the struct contained a pointer or had padding
with uninitialized memory); \texttt{bytemuck} makes it a checked, safe operation instead.

\subsection{The reference grid}
\begin{lstlisting}
fn generate_grid_lines(extent: f32, spacing: f32) -> Vec<Vertex> {
    let mut verts = Vec::new();
    let mut x = -extent;
    while x <= extent {
        verts.push(Vertex { position: [x, -extent] });
        verts.push(Vertex { position: [x, extent] });
        x += spacing;
    }
    // ... same for horizontal lines ...
    verts
}
\end{lstlisting}
A plain function generating a \texttt{Vec<Vertex>} of line endpoints (pairs, since the
pipeline is configured with \texttt{PrimitiveTopology::LineList} — every two vertices is
one independent line segment). This exists purely so the camera's pan/zoom has something
visible to confirm against before Phase 2 introduces real scene elements — it is
explicitly \emph{not} the beginning of the real rendering pipeline (no culling, no LOD, no
caching — all of that is future-phase work per \texttt{PLAN.md}).

\subsection{The shader}
\begin{lstlisting}[language=C]
struct CameraUniform {
    view_proj: mat4x4<f32>,
};

@group(0) @binding(0)
var<uniform> camera: CameraUniform;

@vertex
fn vs_main(in: VertexInput) -> VertexOutput {
    var out: VertexOutput;
    out.clip_position = camera.view_proj * vec4<f32>(in.position, 0.0, 1.0);
    return out;
}

@fragment
fn fs_main(in: VertexOutput) -> @location(0) vec4<f32> {
    return vec4<f32>(0.35, 0.35, 0.4, 1.0);
}
\end{lstlisting}
WGSL is its own language, not Rust, though the syntax rhymes (it was designed around the
same era with similar sensibilities). \texttt{@group(0) @binding(0)} matches the
\texttt{BindGroupLayout} set up in Rust — group 0, binding 0 is where the camera's view
matrix uniform buffer is attached. \texttt{vs\_main} runs once per vertex, transforming
board-space positions into clip space by multiplying by the camera's matrix.
\texttt{fs\_main} runs once per covered pixel, here just returning a flat gray.

\paragraph{The one real bug this phase hit.} The first version of \texttt{fs\_main} was
written as \texttt{fn fs\_main(in: VertexOutput) -> vec4<f32>}, without the
\texttt{@location(0)} attribute on the return type. That compiles fine in Rust (shaders are
just a string embedded via \texttt{include\_str!} until runtime) but \texttt{wgpu} panicked
the moment the app actually ran: \textit{"Entry point arguments and return values must all
have bindings."} WGSL requires every shader input/output to declare which pipeline binding
it corresponds to; a fragment shader's return value implicitly needs \texttt{@location(0)}
to say ``this is color target 0.'' This is exactly why \S\ref{sec:verification} (actually
running the app) matters: \texttt{cargo build} and \texttt{cargo clippy} both passed clean
with this bug present, because shader validation only happens when wgpu actually compiles
the WGSL at runtime.

\subsection{The result enum, not a direct error propagation}
\begin{lstlisting}
pub enum RenderOutcome {
    Rendered,
    Skipped,
    NeedsReconfigure,
}
\end{lstlisting}
Newer wgpu versions changed \texttt{Surface::get\_current\_texture} to return an enum
(\texttt{CurrentSurfaceTexture}) covering several expected-but-not-fatal conditions
(window occluded, surface outdated after a resize, etc.) rather than a
\texttt{Result<\_, SurfaceError>}. Rather than have \texttt{blueboard-app} match on wgpu's
enum directly — which would mean the app crate depending on the exact wgpu version used
inside render — \texttt{Renderer::render} translates it into this small,
render-crate-owned enum. This is the same ``keep the boundary thin'' instinct as the
generic window type above: \texttt{blueboard-app} reacts to three outcomes it actually
cares about, and doesn't need to know wgpu exists.

\section{\texttt{blueboard-app}: the event loop}

File: \texttt{crates/blueboard-app/src/main.rs}.

\subsection{\texttt{ApplicationHandler}, not a closure}
Older \texttt{winit} versions ran an event loop via a closure:
\texttt{event\_loop.run(move |event, elwt| \{ ... \})}. The version resolved by Cargo for
this project (0.30) replaced that with a trait:
\begin{lstlisting}
struct App {
    window: Option<Arc<Window>>,
    renderer: Option<Renderer>,
    camera: Camera,
    dragging: bool,
    last_cursor: DVec2,
}

impl ApplicationHandler for App {
    fn resumed(&mut self, event_loop: &ActiveEventLoop) { /* ... */ }
    fn window_event(&mut self, event_loop: &ActiveEventLoop, window_id: WindowId, event: WindowEvent) { /* ... */ }
}
\end{lstlisting}
\texttt{App} is a plain struct holding all the state the event loop needs between
callbacks. \texttt{window} and \texttt{renderer} are \texttt{Option}s because they don't
exist yet when the program starts — they're created inside \texttt{resumed}, which winit
guarantees is called once the platform backend is actually ready to create windows (on some
platforms, e.g. mobile, a window cannot be created before this point at all, which is why
winit moved to this model). \texttt{\#[derive(Default)]} on \texttt{App} works here because
every field has a sensible default: \texttt{None}, \texttt{Camera::default()} (pan at
origin, 100\% zoom), \texttt{false}, and \texttt{DVec2::ZERO}.

\subsection{Handling events: pattern matching}
\begin{lstlisting}
let (Some(window), Some(renderer)) = (&self.window, &mut self.renderer) else {
    return;
};

match event {
    WindowEvent::CloseRequested => event_loop.exit(),
    WindowEvent::Resized(size) => {
        renderer.resize((size.width, size.height));
        window.request_redraw();
    }
    WindowEvent::RedrawRequested => match renderer.render(&self.camera) {
        RenderOutcome::Rendered | RenderOutcome::Skipped => {}
        RenderOutcome::NeedsReconfigure => { /* ... */ }
    },
    // ...
    _ => {}
}
\end{lstlisting}
Two Rust features doing a lot of work here:
\begin{itemize}
  \item \textbf{\texttt{let-else}}: \texttt{let (Some(window), Some(renderer)) = (...) else
    \{ return; \}} destructures the two \texttt{Option}s and, if either is \texttt{None}
    (i.e. the window isn't created yet), returns early — avoiding a deeply nested
    \texttt{if let Some(...) = ... \{ if let Some(...) = ... \{ ... \} \}}.
  \item \textbf{\texttt{match} on enums}: \texttt{WindowEvent} is an enum with one variant
    per kind of input/system event. The compiler checks match arms are exhaustive — the
    trailing \texttt{\_ => \{\}} catches every event this phase doesn't care about yet (and
    is required for the match to compile, since otherwise not every \texttt{WindowEvent}
    variant would be handled). The nested \texttt{match renderer.render(...)} does the same
    thing for the \texttt{RenderOutcome} enum described above, with \texttt{|} used to
    group two variants (\texttt{Rendered} and \texttt{Skipped}) that need no action.
\end{itemize}

\subsection{Pan and zoom input}
\begin{lstlisting}
WindowEvent::CursorMoved { position, .. } => {
    let cursor = DVec2::new(position.x, position.y);
    if self.dragging {
        self.camera.pan_by_screen_delta(cursor - self.last_cursor);
        window.request_redraw();
    }
    self.last_cursor = cursor;
}
\end{lstlisting}
\texttt{WindowEvent::CursorMoved \{ position, .. \}} destructures just the \texttt{position}
field out of the event struct, with \texttt{..} meaning ``ignore the rest.'' Panning is
computed as the delta between this move and the last recorded cursor position — which is
why \texttt{last\_cursor} has to be tracked across events rather than recomputed each time.
\texttt{window.request\_redraw()} is called explicitly after any state change that should
cause a new frame — winit does not redraw continuously by default; it waits to be told a
redraw is needed, which keeps the app idle (no CPU/GPU usage) when nothing is happening.

\section{Tools used, and why}

\begin{description}
  \item[Nix flake + direnv] (\texttt{flake.nix}, \texttt{.envrc}) pin the exact Rust
    toolchain and all native runtime libraries (Vulkan loader, Wayland/X11 client
    libraries) so the dev environment is reproducible rather than depending on whatever
    happens to be installed system-wide — standard practice on NixOS, decided in Phase 0.
  \item[\texttt{cargo}] Rust's build tool and package manager; \texttt{cargo add} was used
    throughout this phase to let the dependency resolver pick current, mutually-compatible
    versions of \texttt{winit}, \texttt{wgpu}, \texttt{glam}, and \texttt{bytemuck} rather
    than hand-guessing version numbers.
  \item[\texttt{cargo fmt} / \texttt{cargo clippy}] formatting and linting, enforced as a
    \texttt{lefthook} pre-commit hook (Phase 0) so style issues never land in history.
  \item[\texttt{xdotool} and ImageMagick's \texttt{import}] used one-off (via \texttt{nix
    -shell -p}, not added to the project's own dev shell) to actually drive the running app
    — move the mouse, drag, scroll — and screenshot the result, which is how the shader bug
    above was caught and how pan/zoom were confirmed working rather than just assumed from
    reading the code.
\end{description}

\section{Decisions worth calling out}
\label{sec:verification}

\begin{itemize}
  \item \textbf{Camera in \texttt{core}, not \texttt{render}.} Discussed in \S2.1 — the
    payoff was concrete: two unit tests run in milliseconds with no GPU involved.
  \item \textbf{A render-crate-owned \texttt{RenderOutcome} enum instead of re-exporting
    wgpu's error type.} wgpu's own API for this changed shape within the same minor
    version range explored during this phase; insulating the app crate behind a small,
    stable enum means a future wgpu upgrade only has to be reconciled in one place.
  \item \textbf{A static reference grid, not real scene rendering.} Deliberately minimal:
    enough to see the camera work, nothing that would need to be thrown away once Phase 2
    introduces actual vector strokes, tessellation, and culling.
  \item \textbf{Verification meant actually running the app}, not just a clean
    \texttt{cargo build}. The shader bug in \S3.6 would have shipped invisibly otherwise —
    it is a class of bug (runtime shader validation) that Rust's own compiler cannot catch,
    since WGSL lives in a string until wgpu parses it.
\end{itemize}

\section{What's next}
Phase 2, per \texttt{PLAN.md}: the vector pen. Freehand pointer input gets fit to cubic
Bezier paths (the fitting algorithm is already derived in \texttt{planning/math.tex}),
tessellated into the GPU meshes the renderer now knows how to draw, and stored in
\texttt{core}'s scene graph indexed by the R-tree that viewport culling will use starting
in that phase.

\end{document}
